%%═══════════════════════════════════════════════════════════════════
%% Lecture 18 — Fluid Mechanics: Lagrangian vs. Eulerian Formulations
%% 77315 — Computational Physics
%%═══════════════════════════════════════════════════════════════════

\section{Lecture 18 — Fluid Mechanics: Lagrangian \& Eulerian Formulations}\label{sec:lec18}
\hrule\vspace{1.5em}

%%──────────────────────────────────────────────────────────────────
\subsection*{Overview}
Fluid mechanics can be formulated in two complementary viewpoints:
the \emph{Lagrangian} frame (follow a fluid parcel) and the
\emph{Eulerian} frame (observe a fixed point in space).
We derive the continuity, momentum, and energy equations in both
frames and connect them via the material derivative.  As an
application, we derive the speed of sound from linearised Euler equations.

%%──────────────────────────────────────────────────────────────────
\subsection{Lagrangian vs.\ Eulerian Descriptions}\label{subsec:lec18:frames}

\subsubsection*{Material derivative}
\begin{equation}
  \frac{D}{Dt} = \frac{\partial}{\partial t} + \mathbf{u}\cdot\nabla.
  \label{eq:lec18:matderiv}
\end{equation}
The Lagrangian time derivative of a field $q$ following the fluid is
$Dq/Dt$; the Eulerian time derivative is $\partial q/\partial t$.

%%──────────────────────────────────────────────────────────────────
\subsection{Conservation Equations (Eulerian Form)}\label{subsec:lec18:euler}

\subsubsection*{Continuity (mass conservation)}
\begin{equation}
  \frac{\partial \rho}{\partial t} + \nabla\cdot(\rho\mathbf{u}) = 0.
  \label{eq:lec18:cont}
\end{equation}

\subsubsection*{Momentum (Euler equation, inviscid)}
\begin{equation}
  \frac{\partial(\rho\mathbf{u})}{\partial t}
  + \nabla\cdot(\rho\mathbf{u}\mathbf{u}) = -\nabla p.
  \label{eq:lec18:momentum}
\end{equation}
Including viscosity adds a $\mu\nabla^2\mathbf{u}$ term (Navier–Stokes).

\subsubsection*{Energy (1st law of thermodynamics)}
\begin{equation}
  \frac{\partial e}{\partial t}
  + \nabla\cdot\!\left[(e+p)\mathbf{u}\right] = 0,
  \label{eq:lec18:energy}
\end{equation}
where $e = \rho\epsilon + \tfrac{1}{2}\rho|\mathbf{u}|^2$
is the total energy density ($\epsilon$ = specific internal energy).

%%──────────────────────────────────────────────────────────────────
\subsection{Lagrangian Form}\label{subsec:lec18:lagrangian}

Following a fluid element, \eqref{eq:lec18:cont}–\eqref{eq:lec18:energy} become

\begin{align}
  \frac{D\rho}{Dt} + \rho\,\nabla\cdot\mathbf{u} &= 0,
  \label{eq:lec18:lag_cont}\\
  \rho\frac{D\mathbf{u}}{Dt} &= -\nabla p,
  \label{eq:lec18:lag_mom}\\
  \rho\frac{D\epsilon}{Dt} &= -p\,\nabla\cdot\mathbf{u}.
  \label{eq:lec18:lag_energy}
\end{align}

Equation~\eqref{eq:lec18:lag_energy} is the work done by pressure
in compressing the parcel — identical to the thermodynamic identity
$dU = -p\, dV$.

%%──────────────────────────────────────────────────────────────────
\subsection{Sound Wave Derivation}\label{subsec:lec18:sound}

Linearise about a uniform rest state $\rho_0$, $p_0$, $\mathbf{u}=0$:
$\rho = \rho_0+\rho'$, $p=p_0+p'$, small perturbations.

From continuity and momentum:
\begin{align}
  \frac{\partial \rho'}{\partial t} + \rho_0\,\nabla\cdot\mathbf{u}' &= 0, \\
  \rho_0\frac{\partial\mathbf{u}'}{\partial t} &= -\nabla p'.
\end{align}
With an adiabatic EOS $p=p(\rho)$, so $p' = c_s^2\rho'$, eliminate
$\mathbf{u}'$ to get the wave equation:
\begin{equation}
  \frac{\partial^2 \rho'}{\partial t^2} = c_s^2\,\nabla^2\rho',
  \quad
  c_s^2 = \left.\frac{\partial p}{\partial \rho}\right|_s
         = \frac{\gamma p_0}{\rho_0}.
  \label{eq:lec18:sound}
\end{equation}

\nt{The speed of sound $c_s = \sqrt{\gamma p_0/\rho_0}$ sets the CFL
time-step limit for hydrodynamic simulations:
$\Delta t \le \Delta x / c_s$.}

%%──────────────────────────────────────────────────────────────────
\subsection{Comparison of Numerical Frameworks}\label{subsec:lec18:numerics}

\begin{center}
\begin{tabular}{lll}
\toprule
Aspect & Lagrangian & Eulerian \\
\midrule
Grid & moves with fluid & fixed in space \\
Advection & absent (built in) & explicit advection term \\
Interface tracking & exact (trivial) & difficult (VOF, level-set) \\
Large deformations & mesh tangling & no mesh issue \\
\bottomrule
\end{tabular}
\end{center}

\nt{Most astrophysical codes (e.g.\ supernova 1D codes) use the Lagrangian formulation because the fluid compresses dramatically; most engineering CFD codes use Eulerian (finite-volume) formulations.}
